\begin{trapbox}
\textbf{\textcolor{CTrap}{新增提醒：竖直坐标差不是一般的点线距离}}

当 $B\ne0$ 时，$PR=|Ax_0+By_0+C|/|B|$ 是竖直长度。直线不是水平线、且点不在线上时，它大于实际垂直距离；须再乘上 $|B|/\sqrt{A^2+B^2}$。点在直线上时，两者均为零。
这里的 $|B|$ 只来自求同横坐标点，最终距离公式仍适用于 $B=0$ 的竖直直线，不可把中间步骤的限制当成最终公式的限制。
\medskip


	\begin{description}[style=nextline, leftmargin=2.8em, labelsep=0.5em, itemsep=0.35em]
		\item[\textbf{\textcolor{CTrap}{易错点一（分母漏开方）}}]
			没有将分母开平方，直接使用 $A^2+B^2$ 作为分母。
		\item[\textbf{\textcolor{CTrap}{正确理解（分母含平方根）：}}]
			公式分母必须是 $\sqrt{A^2+B^2}$，不是 $A^2+B^2$。
		\item[\textbf{\textcolor{CTrap}{错因分析（记忆混淆）：}}]
			记错为 $d = \frac{|Ax_0+By_0+C|}{{A^2+B^2}}$，忽略开方步骤。
	\end{description}

	\medskip
	\begin{description}[style=nextline, leftmargin=2.8em, labelsep=0.5em, itemsep=0.35em]
		\item[\textbf{\textcolor{CTrap}{易错点二（直线未化一般式）}}]
			将非一般式的直线方程直接代入公式。
		\item[\textbf{\textcolor{CTrap}{正确理解（先化一般式）：}}]
			使用公式前必须将直线方程化为 $Ax+By+C=0$ 的形式。
		\item[\textbf{\textcolor{CTrap}{错因分析（直接代入系数）：}}]
			例如直线移项后为 $2x-y+1=0$，应取 $A=2,B=-1,C=1$；不能漏掉 $y$ 的负号。
	\end{description}

	\medskip
	\begin{description}[style=nextline, leftmargin=2.8em, labelsep=0.5em, itemsep=0.35em]
		\item[\textbf{\textcolor{CTrap}{易错点三（距离为负理解）}}]
			将距离计算结果写为负数。
		\item[\textbf{\textcolor{CTrap}{正确理解（距离非负）：}}]
			距离是长度，必须非负，分子必须加绝对值。
		\item[\textbf{\textcolor{CTrap}{错因分析（忽略绝对值）：}}]
			计算时去掉绝对值符号，可能得到负值，误以为距离可为负。
	\end{description}

\end{trapbox}
